% !TEX root = ../main.tex
% !TEX program = xelatex
\section{模型假设与符号说明}

\subsection{模型假设}

为便于分析与建模，作如下假设：

\noindent\textbf{假设1：}光伏优先满足负载，外网购电可直接供负载或用于充电；富余电量在容量和充电功率允许时存入储能。

\noindent\textbf{假设2：}每个10分钟时段内，负载功率、光伏功率和电价保持不变。

\noindent\textbf{假设3：}问题二至四的储电量跨日连续衔接，日初电量取前一日实际日末电量。

\noindent\textbf{假设4：}决策仅使用当时已发生的观测和已发布的预报，不提前使用未来真实供需及电价。

\noindent\textbf{假设5：}同一时段的调整费用，仅按最终订单相对零点原计划的净变化结算一次。

\noindent\textbf{假设6：}不考虑电池自放电与寿命折旧，未利用的富余光伏不计售电收益。

\newpage
\subsection{符号说明}

设 $d$ 表示日期，$k$ 表示一天中的第 $k$ 个10分钟时段，其中

\begin{equation}
k=1,2,\ldots,144,
\qquad
\Delta t=\frac{1}{6}\ \mathrm{h}.
\end{equation}

$S_{d,0}$ 表示第 $d$ 天0:00的储电量，$S_{d,k}$ 表示第 $k$ 个时段结束后的储电量。问题一仅研究一个代表日，因此相应公式中可省略日期下标 $d$。

主要公共符号如表所示。各问中特有的预测误差、准备量、订单调整量及价格模型参数，在对应模型中首次出现时再单独定义。

\begin{table}[H]
    \centering
    \caption{主要符号及其含义}
    \begin{tabular}{ccl}
        \hline
        符号 & 含义 & 单位或取值 \\
        \hline
        $d$ & 日期下标 & -- \\
        $k$ & 日内10分钟时段下标 & $1,\ldots,144$ \\
        $\Delta t$ & 单个调度时段长度 & $1/6\,\mathrm{h}$ \\
        $L_{d,k}$ & 小区负载功率 & kW \\
        $P^{PV}_{d,k}$ & 光伏发电功率 & kW \\
        $E^L_{d,k}$ & 第 $k$ 时段负载电量 & kWh \\
        $E^{PV}_{d,k}$ & 第 $k$ 时段光伏电量 & kWh \\
        $p_{d,k}$ & 对应时段外网电价 & 元/kWh \\
        $m_{d,k}$ & 光伏直接供给负载的电量 & kWh \\
        $n_{d,k}$ & 光伏直供后的负载缺口 & kWh \\
        $r_{d,k}$ & 光伏直供后的富余电量 & kWh \\
        $q_{d,k}$ & 普通外网购电量 & kWh \\
        $x_{d,k}$ & 外网购电中用于储能充电的电量 & kWh \\
        $y_{d,k}$ & 富余光伏用于储能充电的电量 & kWh \\
        $z_{d,k}$ & 储能设备向负载输出的电量 & kWh \\
        $S_{d,k}$ & 第 $k$ 时段结束后的内部储电量 & kWh \\
        $w_{d,k}$ & 未被负载或储能利用的光伏电量 & kWh \\
        $S_{\min}$ & 储电量运行下限 & 1200 kWh \\
        $S_{\max}$ & 储电量运行上限 & 10800 kWh \\
        $P_{\max}$ & 最大充、放电功率 & 5000 kW \\
        $\eta_c$ & 充电效率 & 0.90 \\
        $\eta_d$ & 放电效率 & 0.90 \\
        \hline
    \end{tabular}
\end{table}

上述符号构成四个问题共用的电量流和储能状态基础，其中设备容量、功率和效率均直接采用题目给定参数，而非额外假设。
